% =============================================================================
% CHAPTER 5 — RESULTS AND DISCUSSION
% =============================================================================
\chapter{Results and Discussion}
\label{ch:results}

\section{System Testing}

The StockSense AI prototype was subjected to a structured testing process covering functional
correctness, boundary condition handling, pipeline integration, and performance measurement.
Testing was conducted using a synthesised dataset comprising ten distinct products mapped to
four suppliers, chosen to exercise a diverse range of inventory states simultaneously — some
products well-stocked, some approaching threshold, and some already at stockout. All test
executions were performed with the n8n instance running locally and connected to the
Supabase cloud database.

Testing was conducted at three levels:

\textbf{Unit-level testing} verified the correctness of individual formula implementations
by comparing computed outputs against hand-calculated expected values for the stockout
prediction, risk classification, and supplier WSM scoring formulas.

\textbf{Workflow-level testing} executed each n8n workflow individually via manual trigger and
verified that the database state transitioned correctly. Each workflow was tested with valid
inputs, missing field inputs, and edge-case inputs (zero stock, zero sales velocity, zero
suppliers).

\textbf{End-to-end pipeline testing} triggered WF-08 and observed the sequential execution of
all sub-workflows, verifying that the final database state (updated stockout dates, fresh alerts,
new suggestions, recalculated supplier scores, saved snapshot, system log entry) was consistent
with expectations.

\section{Test Cases}

Table~\ref{tab:testcases} presents a representative set of test cases drawn from the structured
testing exercise.

{\singlespacing
\renewcommand{\arraystretch}{1.5}
\begin{longtable}{@{} p{1.2cm} >{\raggedright\arraybackslash}p{3.5cm} >{\raggedright\arraybackslash}p{4.2cm} >{\raggedright\arraybackslash}p{4.5cm} p{1.5cm} @{}}
\caption{Test Case Results}
\label{tab:testcases} \\
\toprule
\textbf{TC\#} & \textbf{Test Scenario} & \textbf{Input / Precondition} &
\textbf{Expected Output} & \textbf{Result} \\
\midrule
\endfirsthead
\toprule
\textbf{TC\#} & \textbf{Test Scenario} & \textbf{Input / Precondition} &
\textbf{Expected Output} & \textbf{Result} \\
\midrule
\endhead
\bottomrule
\endfoot

TC01 & Product Add via WF-01 &
  POST with new product\_id, name, stock=50, sales=5 &
  Product row created in DB; webhook returns 200 OK &
  \textbf{PASS} \\

TC02 & Product Update via WF-01 &
  POST with existing product\_id, updated stock=75 &
  Existing product row updated; no duplicate row &
  \textbf{PASS} \\

TC03 & Stockout Calculation — Normal &
  stock=40, avg\_daily\_sales=5 &
  days\_to\_stockout=8; date = today + 8 days &
  \textbf{PASS} \\

TC04 & Stockout Calculation — Zero Sales &
  stock=20, avg\_daily\_sales=0 &
  days\_to\_stockout=999; estimated\_date=NULL &
  \textbf{PASS} \\

TC05 & Stockout Calculation — Zero Stock &
  stock=0, avg\_daily\_sales=3 &
  days\_to\_stockout=0; date=today (immediate) &
  \textbf{PASS} \\

TC06 & Risk Classification — HIGH &
  days\_to\_stockout=2 &
  risk\_level = ``HIGH'' &
  \textbf{PASS} \\

TC07 & Risk Classification — MEDIUM &
  days\_to\_stockout=5 &
  risk\_level = ``MEDIUM'' &
  \textbf{PASS} \\

TC08 & Risk Classification — LOW &
  days\_to\_stockout=15 &
  risk\_level = ``LOW'' &
  \textbf{PASS} \\

TC09 & Alert Generation — Below Threshold &
  stock=10, reorder\_threshold=20 &
  Stock Alert record created with status=ACTIVE &
  \textbf{PASS} \\

TC10 & Alert Generation — Above Threshold &
  stock=30, reorder\_threshold=20 &
  No alert created for this product &
  \textbf{PASS} \\

TC11 & AI Reorder Suggestion — Valid &
  3 HIGH-risk products, 2 active suppliers &
  3 Reorder Suggestion rows with supplier\_id, qty, reason &
  \textbf{PASS} \\

TC12 & AI Reorder — No At-Risk Products &
  All products at LOW risk &
  Empty suggestions array; no DB insert &
  \textbf{PASS} \\

TC13 & Supplier WSM Scoring &
  4 suppliers with varying R, P, delivery\_time &
  Composite scores and grades computed; Suppliers table updated &
  \textbf{PASS} \\

TC14 & WhatsApp Alert Delivery &
  Active alerts present; Twilio sandbox configured &
  Message delivered to registered number within 5 seconds &
  \textbf{PASS} \\

TC15 & WhatsApp Anti-Spam Filter &
  last\_alerted\_at = 2 hours ago for product X &
  Product X excluded from message; only other products alerted &
  \textbf{PASS} \\

TC16 & AI Vision — Product Photo Scan &
  Base64-encoded image of a biscuit packet &
  product\_name and category extracted; unique product\_id returned &
  \textbf{PASS} \\

TC17 & Invoice Smart Scan &
  Image of supplier bill with 3 line items &
  All 3 products fuzzy-matched; new stock levels computed &
  \textbf{PASS} \\

TC18 & Full Pipeline — WF-08 Execution &
  10 products, 4 suppliers; manual trigger &
  All sub-workflows complete; snapshot saved; system log entry written &
  \textbf{PASS} \\

TC19 & Health Score — All Well-Stocked &
  10 products, all LOW risk, none OOS &
  Health score = 100 &
  \textbf{PASS} \\

TC20 & Health Score — Severe Scenario &
  0 LOW, 2 MEDIUM, 0 HIGH, 8 OOS products &
  Health score = 18 &
  \textbf{PASS} \\

\end{longtable}
}

All twenty test cases passed, confirming functional correctness across normal, edge-case, and
integration scenarios.

\section{Screenshots Description}

The following descriptions correspond to screenshot placeholders that should be inserted when
preparing the final printed copy of this report. Screenshots were captured from the live
prototype running against the Supabase cloud database.

\begin{figure}[H]
\centering
\includegraphics[width=0.9\textwidth]{images/ss_dashboard.png}
\caption{StockSense AI --- Main Dashboard with Health Score and KPI Cards.}
\label{fig:ss_dashboard}
\end{figure}

\begin{figure}[H]
\centering
\includegraphics[width=0.9\textwidth]{images/ss_products.png}
\caption{Products Inventory View with Color-Coded Risk Levels.}
\label{fig:ss_products}
\end{figure}

\begin{figure}[H]
\centering
\includegraphics[width=0.9\textwidth]{images/ss_reorder.png}
\caption{AI-Powered Reorder Suggestions based on Weighted Scoring.}
\label{fig:ss_reorder}
\end{figure}

\begin{figure}[H]
\centering
\includegraphics[width=0.9\textwidth]{images/ss_suppliers.png}
\caption{StockSense AI — Supplier Analytics and Scoring Page.}
\label{fig:ss_suppliers}
\end{figure}

\begin{figure}[H]
\centering
\includegraphics[width=0.9\textwidth]{images/ss_pipeline.png}
\caption{StockSense AI — Pipeline Execution with Animated Overlay.}
\label{fig:ss_pipeline}
\end{figure}

\begin{figure}[H]
\centering
\includegraphics[width=0.6\textwidth]{images/ss_whatsapp.png}
\caption{StockSense AI — WhatsApp Stock Alert Message Delivered via Twilio.}
\label{fig:ss_whatsapp}
\end{figure}

\section{Performance Evaluation}

\subsection{Pipeline Execution Time}

The complete end-to-end pipeline (WF-08 orchestrating WF-02 through WF-07) was measured across
five manual trigger executions. The average total execution time was approximately 13 seconds
for a dataset of 10 products and 4 suppliers. The majority of this time is accounted for by
the Groq API call in WF-05 (approximately 4--6 seconds for Llama~3.1~8B inference) and
sequential Supabase write operations across workflows. WhatsApp message delivery via Twilio
was consistently achieved within 5 seconds of the WF-07 execution.

\subsection{Health Score Validation}

Table~\ref{tab:healthscores} validates the health score formula across four distinct inventory
scenarios designed to represent real-world kirana store conditions.

\begin{table}[H]
\centering
\caption{Inventory Health Score Validation Across Four Scenarios}
\label{tab:healthscores}
{\singlespacing
\begin{tabular}{@{}l c c c c c c@{}}
\toprule
\textbf{Scenario} & $n_L$ & $n_M$ & $n_H$ & $n_O$ & $N$ & $H$ \\
\midrule
All well-stocked  & 10 & 0 & 0 & 0  & 10 & 100 \\
Demo state        & 3  & 1 & 6 & 0  & 10 & 55  \\
Severe stockout   & 0  & 2 & 0 & 8  & 10 & 18  \\
Complete stockout & 0  & 0 & 0 & 10 & 10 & 5   \\
\bottomrule
\end{tabular}
}
\end{table}

The formula behaves consistently with real-world intuition: a fully stocked store scores
100; a store with 6 HIGH-risk and 3 LOW-risk products scores 55, correctly indicating a
concerning but not critical state; a store with 8 OOS products scores 18, signalling the
need for urgent intervention; and a completely out-of-stock store scores 5 (not zero, due
to the $n_O \times 5$ weight, which preserves the score's informativeness at the extremes).

\subsection{Comparative Performance Against Existing Solutions}

\begin{table}[H]
\centering
\caption{Comparison of StockSense AI with Existing Inventory Solutions}
\label{tab:comparison}
{\singlespacing
\renewcommand{\arraystretch}{1.3}
\begin{tabularx}{\textwidth}{@{} >{\raggedright\arraybackslash}p{4.8cm} >{\raggedright\arraybackslash}X >{\raggedright\arraybackslash}X >{\raggedright\arraybackslash}X >{\raggedright\arraybackslash}X @{}}
\toprule
\textbf{Feature} & \textbf{StockSense AI} & \textbf{Vyapar} & \textbf{Zoho} & \textbf{Khatabook} \\
\midrule
AI Reorder & \textbf{Yes} & No & Limited & No \\
Stockout Prediction & \textbf{Automated} & Threshold alerts only & Basic & No \\
Supplier Scoring & \textbf{WSM-based} & No & Basic & No \\
WhatsApp Alerts & \textbf{Automated} & Transactional only & Transactional only & No \\
Monthly Cost & \textbf{Rs.\ 0--500} & Rs.\ 0--599 & Rs.\ 2999+ & Rs.\ 0 \\
Setup Complexity & \textbf{Low} & Medium & High & Low \\
Real-time Dashboard & \textbf{Yes} & Limited & Yes & Basic \\
\bottomrule
\end{tabularx}
}
\end{table}

\subsection{Infrastructure Cost Analysis}

During prototype development, the system operated entirely on free-tier services: Supabase
free tier (500~MB database, 50,000 monthly active users), Groq free tier (generous rate
limits for Llama~3.1~8B), and n8n self-hosted on the developer's machine. For production
deployment where n8n must run continuously in the cloud, the estimated monthly cost breakdown
is: n8n hosting on Railway or Render (\rupee{}300--500/month), Twilio WhatsApp per-message
charges (nominal for typical alert volumes), Supabase free tier sufficient for single-store
scale. The total estimated production cost per store thus remains below \rupee{}500 per month,
which is well within the financial reach of a typical kirana store owner.
